\scanpage{19}
- 191 - \hfill n\textsuperscript{o} 391

racines de $R$ sont telles que $(u|u)=1$ ou $(u|u)=2$, (en excluant le cas des systèmes de racines de type $G_2$, dont la structure a été déterminée au § 6, n\textsuperscript{o} 12). Inversement, partons d'un réseau $L$ dans un espace euclidien $V$, et soit $R$ la partie formée de \uline{tous} les vecteurs de $L$ tels que $(u|u)=\lambda$, (ou vérifiant l'une des deux conditions $(u|u)=\lambda$, $(u|u)=2\lambda$), pour un $\lambda$ convenable. Alors, si $R$ engendre $V$ et si on suppose de plus que pour tout couple $(u,v)$ de vecteurs de $R$, le nombre $2(u|v)/(u|u)$ est \uline{entier}, $R$ est un \uline{système de racines} : en effet, si $s_u$ est la réflexion orthogonale de vecteur $u$, $s_u(v)$ appartient alors au sous-groupe de $V$ engendré par $u$ et $v$, donc \uline{a fortiori} à $L$, et en outre $s_u(v)$ a même norme que $v$, donc est dans $R$, ce qui prouve $(\mathrm{SR}_{\mathrm I})$ ; comme $(\mathrm{SR}_{\mathrm{II}})$ et $(\mathrm{SR}_{\mathrm{III}})$ sont vérifiés par hypothèse, cela prouve notre assertion. Il restera naturellement à vérifier dans chaque cas que le système $R$ est réduit et irréductible.

Nous allons voir qu'il est possible d'obtenir tous les systèmes de racines correspondant aux schémas de Dynkin du th. 2, ($G_2$ exclus), en prenant pour $V$ un espace numérique $\mathrm R^n$, dont nous désignerons par $\varepsilon_i$, $(1\leq i\leq n)$, la base canonique, ($(\varepsilon_i^*)$ étant la base duale), et où la forme métrique admet $(\varepsilon_i)$ pour \uline{base orthonormale} ; on considérera alors dans $V$ les réseaux $L_i$, $(i=0,1,2,3)$, suivants notés aussi $L_i(V)$, (et leurs intersections par certains sous-espaces de $V$) :

1) Le réseau $L_{\unsure{o|0}}$ est le $\mathrm Z$-module de base $(\varepsilon_i)$ ; il est clair que $(u|v)$ est entier pour $u$, $v$ dans $L_{\unsure{o|0}}$. En outre, les vecteurs de $L_{\unsure{o|0}}$ tels que $(u|u)=1$ sont les vecteurs $\pm\varepsilon_i$, $(1\leq i\leq n)$ ; ceux pour lesquels $(u|u)=2$ sont les vecteurs $\pm\varepsilon_i\pm\varepsilon_j$ pour $i\neq j$.
\mspage{18}
